\section{图论进阶章正文案例推导补强}

这一节补齐本章正文出现过的 LeetCode 案例推导。阅读顺序统一按“题目说明、样例入口、暴力基线、暴力过程、重复工作、优化条件、相邻状态、状态复用、指针和字段、代码落点、正确性、边界实验、复杂度变化、迁移追问”展开；目的不是新增题量，而是把正文中的案例都讲成可复述、可证明、可迁移的解题链。

\subsection{LeetCode 684 冗余连接}

\textbf{题目说明。} LeetCode 684 冗余连接 的题面入口要先还原成具体任务：树加一条边会形成唯一环。

\textbf{样例入口。} 拿 edges=[[1, 2], [1, 3], [2, 3]] 手算，前两条边把 1、2、3 连成一棵树，第三条边连接的两个点已经同组，所以它就是冗余边。

\textbf{暴力基线。} 慢版本每次合并后用 DFS 或 BFS 重新判断连通性。它正确但重复遍历已有连接；当关系只会合并不会拆开时，可以压成集合代表。放到本题就是：先按“树加一条边会形成唯一环”完整试慢版本；样例里已经能看到“规律是无向图加边时，第一次 union 失败的边形成环”，所以优化前必须先能说清慢版本枚举了哪些候选。

\textbf{暴力过程。} 拿 edges=[[1, 2], [1, 3], [2, 3]] 手算，前两条边把 1、2、3 连成一棵树，第三条边连接的两个点已经同组，所以它就是冗余边。

\textbf{重复工作。} 题面模型：树加一条边会形成唯一环。暴力版的慢点要落到本题：慢版本会围绕“树加一条边会形成唯一环”反复搜索连通性。样例规律是“规律是无向图加边时，第一次 union 失败的边形成环”，说明关系只需要被压成集合代表；每次 union 失败或成功都要能翻译回题意。后面的优化只消掉这类重复工作，不能改变答案集合。

\textbf{优化条件。} 本题的关键观察是：按顺序加边，若一条边两端已连通，则它就是冗余边。这句话必须能翻译成可维护状态；如果题目条件改变后观察不再成立，就不能继续套原来的写法。

\textbf{相邻状态。} 规律是无向图加边时，第一次 union 失败的边形成环。把这个特例继续拆开看：union 只是在维护等价类，不是在保存具体路径。两个节点代表元相同只能说明连通或关系一致；如果题目还要距离、比例或奇偶性，必须把附加信息挂到根关系上。

\textbf{状态复用。} 把连通分量压成代表元；查询只比较代表，合并只发生在两个根之间。本题落点是：规律是无向图加边时，第一次 union 失败的边形成环。手算时只改被新输入影响的那一小部分，而不是回头重算完整候选。

\textbf{指针和字段。} 状态字段要能说清：parent、size 或 rank、find 返回的代表元、合并时的附加信息；带权或奇偶并查集还要维护到根的关系。手算时按这个协议跟踪：拿 edges=[[1, 2], [1, 3], [2, 3]] 手算，前两条边把 1、2、3 连成一棵树，第三条边连接的两个点已经同组，所以它就是冗余边。然后按协议复查：手算时写 parent 表和集合代表。每次 union 只改变根之间的关系，find 后代表相同才说明连通。

\textbf{代码落点。} find 路径压缩不能破坏附加权值；union 只能连接两个根；合并失败和重复边要保留题意解释。关键代码行必须能逐句翻译成“旧状态怎样变成新状态”，否则就只是模板。

\textbf{正确性。} 证明从等价关系开始：合并维护连通性的传递性，find 返回的代表元就是当前集合身份。

\textbf{边界实验。} 先用节点编号、最后一条成环边、重复边、并查集初始化大小做反例检查。实验时覆盖节点编号、最后一条成环边、重复边、并查集初始化大小，打印每个元素的代表元和集合大小。重复合并、已经连通的边和字符串编号最容易暴露实现问题。

\textbf{复杂度变化。} 暴力每次查询都搜索连通性代价高；并查集把合并和查询摊还到近似常数，整体接近线性。

\textbf{迁移追问。} 如果题目改成“有向冗余连接需要同时处理入度为二和环，复杂度更高”，先判断原状态是否仍然覆盖新答案集合，再决定是微调边界、改 value 语义，还是换算法。

\subsection{LeetCode 547 省份数量}

\textbf{题目说明。} LeetCode 547 省份数量 的题面入口要先还原成具体任务：城市连接关系形成无向图，省份就是连通分量。

\textbf{样例入口。} 拿 isConnected 中 0 和 1 相连、2 独立手算，0 和 1 union 后属于同一代表，2 保持自己，省份数为 2。

\textbf{暴力基线。} 慢版本每次合并后用 DFS 或 BFS 重新判断连通性。它正确但重复遍历已有连接；当关系只会合并不会拆开时，可以压成集合代表。放到本题就是：先按“城市连接关系形成无向图，省份就是连通分量”完整试慢版本；样例里已经能看到“规律是矩阵中的一条连接边对应集合合并，最后代表元数量就是连通分量数”，所以优化前必须先能说清慢版本枚举了哪些候选。

\textbf{暴力过程。} 拿 isConnected 中 0 和 1 相连、2 独立手算，0 和 1 union 后属于同一代表，2 保持自己，省份数为 2。

\textbf{重复工作。} 题面模型：城市连接关系形成无向图，省份就是连通分量。暴力版的慢点要落到本题：慢版本会围绕“城市连接关系形成无向图，省份就是连通分量”反复搜索连通性。样例规律是“规律是矩阵中的一条连接边对应集合合并，最后代表元数量就是连通分量数”，说明关系只需要被压成集合代表；每次 union 失败或成功都要能翻译回题意。后面的优化只消掉这类重复工作，不能改变答案集合。

\textbf{优化条件。} 本题的关键观察是：并查集合并所有直接连接的城市，最后统计根数量。这句话必须能翻译成可维护状态；如果题目条件改变后观察不再成立，就不能继续套原来的写法。

\textbf{相邻状态。} 规律是矩阵中的一条连接边对应集合合并，最后代表元数量就是连通分量数。把这个特例继续拆开看：union 只是在维护等价类，不是在保存具体路径。两个节点代表元相同只能说明连通或关系一致；如果题目还要距离、比例或奇偶性，必须把附加信息挂到根关系上。

\textbf{状态复用。} 把连通分量压成代表元；查询只比较代表，合并只发生在两个根之间。本题落点是：规律是矩阵中的一条连接边对应集合合并，最后代表元数量就是连通分量数。手算时只改被新输入影响的那一小部分，而不是回头重算完整候选。

\textbf{指针和字段。} 状态字段要能说清：parent、size 或 rank、find 返回的代表元、合并时的附加信息；带权或奇偶并查集还要维护到根的关系。手算时按这个协议跟踪：拿 isConnected 中 0 和 1 相连、2 独立手算，0 和 1 union 后属于同一代表，2 保持自己，省份数为 2。然后按协议复查：手算时写 parent 表和集合代表。每次 union 只改变根之间的关系，find 后代表相同才说明连通。

\textbf{代码落点。} find 路径压缩不能破坏附加权值；union 只能连接两个根；合并失败和重复边要保留题意解释。关键代码行必须能逐句翻译成“旧状态怎样变成新状态”，否则就只是模板。

\textbf{正确性。} 证明从等价关系开始：合并维护连通性的传递性，find 返回的代表元就是当前集合身份。

\textbf{边界实验。} 先用邻接矩阵对称、自连接、孤立城市、只扫描上三角做反例检查。实验时覆盖邻接矩阵对称、自连接、孤立城市、只扫描上三角，打印每个元素的代表元和集合大小。重复合并、已经连通的边和字符串编号最容易暴露实现问题。

\textbf{复杂度变化。} 暴力每次查询都搜索连通性代价高；并查集把合并和查询摊还到近似常数，整体接近线性。

\textbf{迁移追问。} 如果题目改成“BFS 也可做；并查集更突出动态合并”，先判断原状态是否仍然覆盖新答案集合，再决定是微调边界、改 value 语义，还是换算法。

\subsection{LeetCode 743 网络延迟时间}

\textbf{题目说明。} LeetCode 743 网络延迟时间 的题面入口要先还原成具体任务：有向正权图从起点到所有节点的最短路最大值。

\textbf{样例入口。} 拿 times=[[2, 1, 1], [2, 3, 1], [3, 4, 1]]、起点 2 手算，2 到 1 和 3 距离都是 1，到 4 距离是 2。

\textbf{暴力基线。} 暴力枚举从起点到各点的所有路径，或错误地用普通 BFS 按边数扩展。

\textbf{暴力过程。} 以 times=[[2, 1, 1], [2, 3, 1], [3, 4, 1]]、起点 2 为例，2 到 1 和 3 的距离是 1，到 4 的距离是 2。

\textbf{重复工作。} 题面模型：有向正权图从起点到所有节点的最短路最大值。暴力版的慢点要落到本题：浪费在路径数量可能爆炸，而且普通队列不能保证带权路径的总代价最小。需要每次扩展当前距离最小的候选。后面的优化只消掉这类重复工作，不能改变答案集合。

\textbf{优化条件。} 本题的关键观察是：Dijkstra 求每个节点最短距离，最后取最大，若有无穷则不可达。这句话必须能翻译成可维护状态；如果题目条件改变后观察不再成立，就不能继续套原来的写法。

\textbf{相邻状态。} 规律是非负边最短路用 Dijkstra，最后取所有最短距离的最大值；不可达节点导致答案 -1。把这个特例继续拆开看：队列或堆弹出的顺序必须和代价规则一致；旧状态被跳过，是因为已经有更小代价到达同一状态。只要边权或代价合成方式改变，就要重新证明扩展顺序。

\textbf{状态复用。} Dijkstra：dist[node] 保存当前最短距离，小顶堆保存 (distance, node)。弹出过期距离时跳过，松弛所有出边。手算时只改被新输入影响的那一小部分，而不是回头重算完整候选。

\textbf{指针和字段。} 状态字段要能说清：dist[u] 是已知最短上界，heap top 是当前最小候选，edge weight 是从 u 到 v 的传播时间。手算时按这个协议跟踪：起点 2 距离 0 入堆；弹出 2 后更新 1 和 3 为 1；弹出 3 后更新 4 为 2。

\textbf{代码落点。} 关键是 priority\_queue 默认大顶堆，C++ 要用 greater 或存负距离；弹出时若 d!=dist[u] 就跳过旧状态。关键代码行必须能逐句翻译成“旧状态怎样变成新状态”，否则就只是模板。

\textbf{正确性。} 边权非负时，堆中最小距离节点一旦弹出，任何绕路都不会得到更小距离，因此它的 dist 已最终确定。

\textbf{边界实验。} 先用节点编号从一开始、不可达、重边、过期堆元素做反例检查。实验时覆盖节点编号从一开始、不可达、重边、过期堆元素，打印每次弹出的代价和是否过期。若普通 BFS 与 Dijkstra 结果不同，要能解释边权条件造成的差异。

\textbf{复杂度变化。} 枚举路径指数级；Dijkstra O((V+E)log V) 时间，O(V+E) 空间。

\textbf{迁移追问。} 如果题目改成“边权相同才可用普通 BFS”，先判断原状态是否仍然覆盖新答案集合，再决定是微调边界、改 value 语义，还是换算法。

\subsection{LeetCode 1631 最小体力消耗路径}

\textbf{题目说明。} LeetCode 1631 最小体力消耗路径 的题面入口要先还原成具体任务：路径代价是沿途最大高度差，不是边权和。

\textbf{样例入口。} 拿 heights=[[1, 2, 2], [3, 8, 2], [5, 3, 5]] 手算，路径代价是相邻高度差的最大值，不是差值之和。

\textbf{暴力基线。} 慢版本可以枚举路径，或先用普通队列试跑一个小图。它会暴露步数最少和代价最小是否一致；一旦不一致，就必须围绕代价组织状态。放到本题就是：先按“路径代价是沿途最大高度差，不是边权和”完整试慢版本；样例里已经能看到“规律是 Dijkstra 的距离定义为到达某格的最小最大边权，或二分最大体力做可达性”，所以优化前必须先能说清慢版本枚举了哪些候选。

\textbf{暴力过程。} 拿 heights=[[1, 2, 2], [3, 8, 2], [5, 3, 5]] 手算，路径代价是相邻高度差的最大值，不是差值之和。

\textbf{重复工作。} 题面模型：路径代价是沿途最大高度差，不是边权和。暴力版的慢点要落到本题：慢版本会围绕“路径代价是沿途最大高度差，不是边权和”枚举路径或按错误顺序扩展状态。样例规律是“规律是 Dijkstra 的距离定义为到达某格的最小最大边权，或二分最大体力做可达性”，说明队列或堆的弹出顺序必须和代价定义一致，旧状态为什么能跳过也要讲清。后面的优化只消掉这类重复工作，不能改变答案集合。

\textbf{优化条件。} 本题的关键观察是：Dijkstra 的松弛改成取当前代价和新边代价的最大值。这句话必须能翻译成可维护状态；如果题目条件改变后观察不再成立，就不能继续套原来的写法。

\textbf{相邻状态。} 规律是 Dijkstra 的距离定义为到达某格的最小最大边权，或二分最大体力做可达性。把这个特例继续拆开看：队列或堆弹出的顺序必须和代价规则一致；旧状态被跳过，是因为已经有更小代价到达同一状态。只要边权或代价合成方式改变，就要重新证明扩展顺序。

\textbf{状态复用。} 把路径枚举压成距离数组和优先队列；每次只扩展当前最有希望的未过期状态。本题落点是：规律是 Dijkstra 的距离定义为到达某格的最小最大边权，或二分最大体力做可达性。手算时只改被新输入影响的那一小部分，而不是回头重算完整候选。

\textbf{指针和字段。} 状态字段要能说清：距离数组、优先队列状态、松弛候选、旧状态判定、不可达哨兵；资源约束题还要把资源维度放进状态。手算时按这个协议跟踪：拿 heights=[[1, 2, 2], [3, 8, 2], [5, 3, 5]] 手算，路径代价是相邻高度差的最大值，不是差值之和。然后按协议复查：手算时记录每次弹出的代价、当前距离数组和松弛结果。旧状态被弹出时为什么跳过，也要写在表里。

\textbf{代码落点。} 松弛时只在候选更优时更新；priority\_queue 弹出旧状态要跳过；距离加法用更宽整数。关键代码行必须能逐句翻译成“旧状态怎样变成新状态”，否则就只是模板。

\textbf{正确性。} 证明从扩展顺序开始：当前弹出的状态在代价规则下已经不会被未来路径改得更优。

\textbf{边界实验。} 先用单格、平坦网格、必须绕路、多个相同代价做反例检查。实验时覆盖单格、平坦网格、必须绕路、多个相同代价，打印每次弹出的代价和是否过期。若普通 BFS 与 Dijkstra 结果不同，要能解释边权条件造成的差异。

\textbf{复杂度变化。} 暴力枚举路径可能指数级；Dijkstra 在邻接表和堆实现下按边数、点数和堆操作计算，0-1 BFS 可按节点数加边数计算。

\textbf{迁移追问。} 如果题目改成“也可二分体力上限后 BFS 检查可达”，先判断原状态是否仍然覆盖新答案集合，再决定是微调边界、改 value 语义，还是换算法。

\subsection{LeetCode 127 单词接龙}

\textbf{题目说明。} LeetCode 127 单词接龙 的题面入口要先还原成具体任务：每个单词是节点，一次修改一个字符形成边。

\textbf{样例入口。} 拿 hit -> hot -> dot -> dog -> cog 手算，每次只能改一个字符，所以 BFS 第一层是 hot，第二层才到 dot 或 lot。

\textbf{暴力基线。} 暴力可以把每个单词和所有其他单词比较一遍，差一个字符就连边，再从 beginWord 做 BFS。

\textbf{暴力过程。} 以 hit 到 cog 为例，hit 只能一步到 hot；hot 下一层到 dot 或 lot。第一次到 cog 的层数就是最短转换长度。

\textbf{重复工作。} 题面模型：每个单词是节点，一次修改一个字符形成边。暴力版的慢点要落到本题：浪费在找邻居时每次扫描整个词表。单词长度固定时，可以把 h*t、*ot 这类通配模式作为桶，快速找到只差一位的邻居。后面的优化只消掉这类重复工作，不能改变答案集合。

\textbf{优化条件。} 本题的关键观察是：通配模式桶加速邻居查询，BFS 第一次到达终点就是最短转换长度。这句话必须能翻译成可维护状态；如果题目条件改变后观察不再成立，就不能继续套原来的写法。

\textbf{相邻状态。} 规律是单词作为节点，通配桶如 h*t 把邻居快速找出；第一次到达终点就是最短转换长度。把这个特例继续拆开看：状态节点不一定等于原始位置，还可能包含钥匙、剩余资源、时间或访问集合。visited 只能合并未来能力完全相同的状态，否则会把合法路径剪掉。

\textbf{状态复用。} 预处理 pattern -> words 的映射，BFS 队列保存当前单词和步数，visited 防止重复入队。手算时只改被新输入影响的那一小部分，而不是回头重算完整候选。

\textbf{指针和字段。} 状态字段要能说清：word 是当前单词，pattern 是替换某一位后的通配键，distance 是转换层数。手算时按这个协议跟踪：hit 生成 *it、h*t、hi*，其中 h*t 桶找到 hot；hot 再通过 *ot 找到 dot 和 lot。

\textbf{代码落点。} 关键是访问标记在入队时设置。桶用完后可以清空，避免不同节点重复展开同一批邻居。关键代码行必须能逐句翻译成“旧状态怎样变成新状态”，否则就只是模板。

\textbf{正确性。} 每条转换边权都为 1，BFS 按层扩展；第一次到达 endWord 时经过的边数最少。

\textbf{边界实验。} 先用终点不在词表、起点和终点相同、桶重复扫描、单词长度一致做反例检查。实验时覆盖终点不在词表、起点和终点相同、桶重复扫描、单词长度一致，并打印入队时的完整状态。若同一位置但资源不同的状态被合并，很多图题会出现隐蔽错误。

\textbf{复杂度变化。} 两两建图 O(\ensuremath{N^{2}}L)，通配桶建图和 BFS 约 O(\ensuremath{NL^{2}}) 或按桶规模计，空间 O(NL)。

\textbf{迁移追问。} 如果题目改成“双向 BFS 可以进一步降低搜索空间”，先判断原状态是否仍然覆盖新答案集合，再决定是微调边界、改 value 语义，还是换算法。

\subsection{LeetCode 1091 二进制矩阵中的最短路径}

\textbf{题目说明。} LeetCode 1091 二进制矩阵中的最短路径 的题面入口要先还原成具体任务：八方向网格无权边，答案是从左上到右下的最少格子数。

\textbf{样例入口。} 拿 [[0, 1], [1, 0]] 手算，从左上可以斜着一步到右下，路径长度是 2；若起点或终点为 1 直接不可达。

\textbf{暴力基线。} 慢版本先按题意画出完整状态图，用普通搜索跑通可达性。这个过程用于确认八方向网格无权边，答案是从左上到右下的最少格子数的节点和边，之后再讨论访问标记、层次或代价顺序。放到本题就是：先按“八方向网格无权边，答案是从左上到右下的最少格子数”完整试慢版本；样例里已经能看到“规律是八方向 BFS，第一次到达终点就是最短路径长度”，所以优化前必须先能说清慢版本枚举了哪些候选。

\textbf{暴力过程。} 拿 [[0, 1], [1, 0]] 手算，从左上可以斜着一步到右下，路径长度是 2；若起点或终点为 1 直接不可达。

\textbf{重复工作。} 题面模型：八方向网格无权边，答案是从左上到右下的最少格子数。暴力版的慢点要落到本题：慢版本会围绕“八方向网格无权边，答案是从左上到右下的最少格子数”反复展开同一批状态。样例规律是“规律是八方向 BFS，第一次到达终点就是最短路径长度”，说明必须把节点、边、访问标记和资源维度一次建清；同一状态不应重复入队，不同资源状态也不能误合并。后面的优化只消掉这类重复工作，不能改变答案集合。

\textbf{优化条件。} 本题的关键观察是：BFS 入队时标记访问，层数或距离随状态保存。这句话必须能翻译成可维护状态；如果题目条件改变后观察不再成立，就不能继续套原来的写法。

\textbf{相邻状态。} 规律是八方向 BFS，第一次到达终点就是最短路径长度。把这个特例继续拆开看：状态节点不一定等于原始位置，还可能包含钥匙、剩余资源、时间或访问集合。visited 只能合并未来能力完全相同的状态，否则会把合法路径剪掉。

\textbf{状态复用。} 把重复搜索压成访问标记、距离数组或拓扑状态；邻居生成也要从全量扫描变成结构化枚举。本题落点是：规律是八方向 BFS，第一次到达终点就是最短路径长度。手算时只改被新输入影响的那一小部分，而不是回头重算完整候选。

\textbf{指针和字段。} 状态字段要能说清：状态编号、邻接生成方式、访问标记、距离或层数、队列内容；若状态含资源，visited 不能只按位置记录。手算时按这个协议跟踪：拿 [[0, 1], [1, 0]] 手算，从左上可以斜着一步到右下，路径长度是 2；若起点或终点为 1 直接不可达。然后按协议复查：手算时画队列或栈，状态必须包含题目需要的所有资源。入队时标记还是出队时标记，要和去重语义一致。

\textbf{代码落点。} 入队时标记还是出队时标记必须一致；方向数组集中定义；多源 BFS 要一次性把所有源点放入初始队列。关键代码行必须能逐句翻译成“旧状态怎样变成新状态”，否则就只是模板。

\textbf{正确性。} 证明从可达性开始：搜索覆盖所有合法边能到达的状态，访问标记只删除重复状态而不删除不同未来能力。

\textbf{边界实验。} 先用起点或终点阻塞、单格、八方向越界、无路可达做反例检查。实验时覆盖起点或终点阻塞、单格、八方向越界、无路可达，并打印入队时的完整状态。若同一位置但资源不同的状态被合并，很多图题会出现隐蔽错误。

\textbf{复杂度变化。} 暴力重复从多个起点搜索会反复遍历图；正确建模和访问标记后，BFS 或 DFS 通常按节点数加边数计算，状态图题则按完整状态数计算。

\textbf{迁移追问。} 如果题目改成“边权相同才能用 BFS，若每格有代价就要最短路”，先判断原状态是否仍然覆盖新答案集合，再决定是微调边界、改 value 语义，还是换算法。

\subsection{LeetCode 1368 使网格图至少有一条有效路径的最小代价}

\textbf{题目说明。} LeetCode 1368 使网格图至少有一条有效路径的最小代价 的题面入口要先还原成具体任务：沿着格子箭头走代价为零，改方向走代价为一。

\textbf{样例入口。} 拿网格箭头指向右和下手算，沿箭头走的代价是 0，改方向才付 1。

\textbf{暴力基线。} 慢版本可以枚举路径，或先用普通队列试跑一个小图。它会暴露步数最少和代价最小是否一致；一旦不一致，就必须围绕代价组织状态。放到本题就是：先按“沿着格子箭头走代价为零，改方向走代价为一”完整试慢版本；样例里已经能看到“规律是边权只有 0 和 1，可用 0-1 BFS；状态是格子，松弛代价取决于当前箭头是否指向邻居”，所以优化前必须先能说清慢版本枚举了哪些候选。

\textbf{暴力过程。} 拿网格箭头指向右和下手算，沿箭头走的代价是 0，改方向才付 1。

\textbf{重复工作。} 题面模型：沿着格子箭头走代价为零，改方向走代价为一。暴力版的慢点要落到本题：慢版本会围绕“沿着格子箭头走代价为零，改方向走代价为一”枚举路径或按错误顺序扩展状态。样例规律是“规律是边权只有 0 和 1，可用 0-1 BFS；状态是格子，松弛代价取决于当前箭头是否指向邻居”，说明队列或堆的弹出顺序必须和代价定义一致，旧状态为什么能跳过也要讲清。后面的优化只消掉这类重复工作，不能改变答案集合。

\textbf{优化条件。} 本题的关键观察是：边权只有零和一，用 0-1 BFS：零代价邻居入队首，一代价邻居入队尾。这句话必须能翻译成可维护状态；如果题目条件改变后观察不再成立，就不能继续套原来的写法。

\textbf{相邻状态。} 规律是边权只有 0 和 1，可用 0-1 BFS；状态是格子，松弛代价取决于当前箭头是否指向邻居。把这个特例继续拆开看：队列或堆弹出的顺序必须和代价规则一致；旧状态被跳过，是因为已经有更小代价到达同一状态。只要边权或代价合成方式改变，就要重新证明扩展顺序。

\textbf{状态复用。} 把路径枚举压成距离数组和优先队列；每次只扩展当前最有希望的未过期状态。本题落点是：规律是边权只有 0 和 1，可用 0-1 BFS；状态是格子，松弛代价取决于当前箭头是否指向邻居。手算时只改被新输入影响的那一小部分，而不是回头重算完整候选。

\textbf{指针和字段。} 状态字段要能说清：距离数组、优先队列状态、松弛候选、旧状态判定、不可达哨兵；资源约束题还要把资源维度放进状态。手算时按这个协议跟踪：拿网格箭头指向右和下手算，沿箭头走的代价是 0，改方向才付 1。然后按协议复查：手算时记录每次弹出的代价、当前距离数组和松弛结果。旧状态被弹出时为什么跳过，也要写在表里。

\textbf{代码落点。} 松弛时只在候选更优时更新；priority\_queue 弹出旧状态要跳过；距离加法用更宽整数。关键代码行必须能逐句翻译成“旧状态怎样变成新状态”，否则就只是模板。

\textbf{正确性。} 证明从扩展顺序开始：当前弹出的状态在代价规则下已经不会被未来路径改得更优。

\textbf{边界实验。} 先用单格、箭头指向越界、多个零代价路径、过期状态做反例检查。实验时覆盖单格、箭头指向越界、多个零代价路径、过期状态，打印每次弹出的代价和是否过期。若普通 BFS 与 Dijkstra 结果不同，要能解释边权条件造成的差异。

\textbf{复杂度变化。} 暴力枚举路径可能指数级；Dijkstra 在邻接表和堆实现下按边数、点数和堆操作计算，0-1 BFS 可按节点数加边数计算。

\textbf{迁移追问。} 如果题目改成“普通 BFS 按步数扩展，不等于总代价最小”，先判断原状态是否仍然覆盖新答案集合，再决定是微调边界、改 value 语义，还是换算法。

\subsection{LeetCode 787 K 站中转内最便宜的航班}

\textbf{题目说明。} LeetCode 787 K 站中转内最便宜的航班 的题面入口要先还原成具体任务：路径代价受中转次数限制，不能只按费用最小弹出节点后永久确定。

\textbf{样例入口。} 拿航班 0->1 价格 100、1->2 价格 100、0->2 价格 500，K=1 手算，允许一次中转时 0->1->2 更便宜。

\textbf{暴力基线。} 慢版本可以枚举路径，或先用普通队列试跑一个小图。它会暴露步数最少和代价最小是否一致；一旦不一致，就必须围绕代价组织状态。放到本题就是：先按“路径代价受中转次数限制，不能只按费用最小弹出节点后永久确定”完整试慢版本；样例里已经能看到“规律是状态必须包含已用中转次数或边数，不能只保存每个城市一个最低价格”，所以优化前必须先能说清慢版本枚举了哪些候选。

\textbf{暴力过程。} 拿航班 0->1 价格 100、1->2 价格 100、0->2 价格 500，K=1 手算，允许一次中转时 0->1->2 更便宜。

\textbf{重复工作。} 题面模型：路径代价受中转次数限制，不能只按费用最小弹出节点后永久确定。暴力版的慢点要落到本题：慢版本会围绕“路径代价受中转次数限制，不能只按费用最小弹出节点后永久确定”枚举路径或按错误顺序扩展状态。样例规律是“规律是状态必须包含已用中转次数或边数，不能只保存每个城市一个最低价格”，说明队列或堆的弹出顺序必须和代价定义一致，旧状态为什么能跳过也要讲清。后面的优化只消掉这类重复工作，不能改变答案集合。

\textbf{优化条件。} 本题的关键观察是：用按边数分层的 Bellman-Ford 或带层数状态的优先队列，状态包含城市和已用边数。这句话必须能翻译成可维护状态；如果题目条件改变后观察不再成立，就不能继续套原来的写法。

\textbf{相邻状态。} 规律是状态必须包含已用中转次数或边数，不能只保存每个城市一个最低价格。把这个特例继续拆开看：队列或堆弹出的顺序必须和代价规则一致；旧状态被跳过，是因为已经有更小代价到达同一状态。只要边权或代价合成方式改变，就要重新证明扩展顺序。

\textbf{状态复用。} 把路径枚举压成距离数组和优先队列；每次只扩展当前最有希望的未过期状态。本题落点是：规律是状态必须包含已用中转次数或边数，不能只保存每个城市一个最低价格。手算时只改被新输入影响的那一小部分，而不是回头重算完整候选。

\textbf{指针和字段。} 状态字段要能说清：距离数组、优先队列状态、松弛候选、旧状态判定、不可达哨兵；资源约束题还要把资源维度放进状态。手算时按这个协议跟踪：拿航班 0->1 价格 100、1->2 价格 100、0->2 价格 500，K=1 手算，允许一次中转时 0->1->2 更便宜。然后按协议复查：手算时记录每次弹出的代价、当前距离数组和松弛结果。旧状态被弹出时为什么跳过，也要写在表里。

\textbf{代码落点。} 松弛时只在候选更优时更新；priority\_queue 弹出旧状态要跳过；距离加法用更宽整数。关键代码行必须能逐句翻译成“旧状态怎样变成新状态”，否则就只是模板。

\textbf{正确性。} 证明从扩展顺序开始：当前弹出的状态在代价规则下已经不会被未来路径改得更优。

\textbf{边界实验。} 先用起点终点相同、不可达、K 为零、便宜路径中转过多做反例检查。实验时覆盖起点终点相同、不可达、K 为零、便宜路径中转过多，打印每次弹出的代价和是否过期。若普通 BFS 与 Dijkstra 结果不同，要能解释边权条件造成的差异。

\textbf{复杂度变化。} 暴力枚举路径可能指数级；Dijkstra 在邻接表和堆实现下按边数、点数和堆操作计算，0-1 BFS 可按节点数加边数计算。

\textbf{迁移追问。} 如果题目改成“若去掉中转限制，普通 Dijkstra 才能直接按费用组织扩展”，先判断原状态是否仍然覆盖新答案集合，再决定是微调边界、改 value 语义，还是换算法。

\subsection{LeetCode 1462 课程表 IV}

\textbf{题目说明。} LeetCode 1462 课程表 IV 的题面入口要先还原成具体任务：大量先修关系查询需要预处理课程之间的可达性。

\textbf{样例入口。} 拿 0->1->2 和查询 0 是否先修 2 手算，直接边只能回答 0->1、1->2，不能回答传递关系。若查询很多，每次 DFS 会重复走相同路径。

\textbf{暴力基线。} 慢版本先按题意画出完整状态图，用普通搜索跑通可达性。这个过程用于确认大量先修关系查询需要预处理课程之间的可达性的节点和边，之后再讨论访问标记、层次或代价顺序。放到本题就是：先按“大量先修关系查询需要预处理课程之间的可达性”完整试慢版本；样例里已经能看到“规律是先做可达性闭包或拓扑合并祖先集合，再 O(1) 回答查询”，所以优化前必须先能说清慢版本枚举了哪些候选。

\textbf{暴力过程。} 拿 0->1->2 和查询 0 是否先修 2 手算，直接边只能回答 0->1、1->2，不能回答传递关系。若查询很多，每次 DFS 会重复走相同路径。

\textbf{重复工作。} 题面模型：大量先修关系查询需要预处理课程之间的可达性。暴力版的慢点要落到本题：慢版本会围绕“大量先修关系查询需要预处理课程之间的可达性”反复展开同一批状态。样例规律是“规律是先做可达性闭包或拓扑合并祖先集合，再 O(1) 回答查询”，说明必须把节点、边、访问标记和资源维度一次建清；同一状态不应重复入队，不同资源状态也不能误合并。后面的优化只消掉这类重复工作，不能改变答案集合。

\textbf{优化条件。} 本题的关键观察是：可以用 Floyd 传递闭包、每个点 BFS，或拓扑 DP 合并祖先集合，查询时直接查可达表。这句话必须能翻译成可维护状态；如果题目条件改变后观察不再成立，就不能继续套原来的写法。

\textbf{相邻状态。} 规律是先做可达性闭包或拓扑合并祖先集合，再 O(1) 回答查询。把这个特例继续拆开看：状态节点不一定等于原始位置，还可能包含钥匙、剩余资源、时间或访问集合。visited 只能合并未来能力完全相同的状态，否则会把合法路径剪掉。

\textbf{状态复用。} 把重复搜索压成访问标记、距离数组或拓扑状态；邻居生成也要从全量扫描变成结构化枚举。本题落点是：规律是先做可达性闭包或拓扑合并祖先集合，再 O(1) 回答查询。手算时只改被新输入影响的那一小部分，而不是回头重算完整候选。

\textbf{指针和字段。} 状态字段要能说清：状态编号、邻接生成方式、访问标记、距离或层数、队列内容；若状态含资源，visited 不能只按位置记录。手算时按这个协议跟踪：拿 0->1->2 和查询 0 是否先修 2 手算，直接边只能回答 0->1、1->2，不能回答传递关系。若查询很多，每次 DFS 会重复走相同路径。然后按协议复查：手算时画队列或栈，状态必须包含题目需要的所有资源。入队时标记还是出队时标记，要和去重语义一致。

\textbf{代码落点。} 入队时标记还是出队时标记必须一致；方向数组集中定义；多源 BFS 要一次性把所有源点放入初始队列。关键代码行必须能逐句翻译成“旧状态怎样变成新状态”，否则就只是模板。

\textbf{正确性。} 证明从可达性开始：搜索覆盖所有合法边能到达的状态，访问标记只删除重复状态而不删除不同未来能力。

\textbf{边界实验。} 先用课程之间没有关系、自身是否算先修、查询很多、图中按题意应为 DAG做反例检查。实验时覆盖课程之间没有关系、自身是否算先修、查询很多、图中按题意应为 DAG，并打印入队时的完整状态。若同一位置但资源不同的状态被合并，很多图题会出现隐蔽错误。

\textbf{复杂度变化。} 暴力重复从多个起点搜索会反复遍历图；正确建模和访问标记后，BFS 或 DFS 通常按节点数加边数计算，状态图题则按完整状态数计算。

\textbf{迁移追问。} 如果题目改成“如果查询很少，逐次搜索可能更省预处理成本；多查询时闭包更稳”，先判断原状态是否仍然覆盖新答案集合，再决定是微调边界、改 value 语义，还是换算法。

\subsection{LeetCode 207 课程表}

\textbf{题目说明。} LeetCode 207 课程表 的题面入口要先还原成具体任务：课程是节点，先修关系是有向边，问题是判断是否有环。

\textbf{样例入口。} 拿课程 0->1 手算，入度为 0 的课程 0 先入队，学完后 1 的入度降为 0；若有 0->1 和 1->0，队列一开始为空。

\textbf{暴力基线。} 暴力可以不断扫描所有课程，找当前先修都已满足的课程；若一轮没有新课程可学，就判断有环。

\textbf{暴力过程。} 以 0->1 为例，课程 0 入度为 0，先学 0，随后 1 的入度变 0。若有 0->1 和 1->0，一开始没有入度为 0 的课。

\textbf{重复工作。} 题面模型：课程是节点，先修关系是有向边，问题是判断是否有环。暴力版的慢点要落到本题：浪费在每轮都重新检查所有课程的先修。入度正好表示还剩多少未满足依赖，学完一门课只影响它指向的后继。后面的优化只消掉这类重复工作，不能改变答案集合。

\textbf{优化条件。} 本题的关键观察是：入度拓扑排序不断学习当前无依赖课程。这句话必须能翻译成可维护状态；如果题目条件改变后观察不再成立，就不能继续套原来的写法。

\textbf{相邻状态。} 规律是拓扑排序用入度为 0 的节点推进，处理数量小于课程数说明有环。把这个特例继续拆开看：状态节点不一定等于原始位置，还可能包含钥匙、剩余资源、时间或访问集合。visited 只能合并未来能力完全相同的状态，否则会把合法路径剪掉。

\textbf{状态复用。} 构建邻接表和 indegree。把所有入度为 0 的课程入队，弹出课程时让后继入度减一，减到 0 就入队。手算时只改被新输入影响的那一小部分，而不是回头重算完整候选。

\textbf{指针和字段。} 状态字段要能说清：indegree[v] 是课程 v 尚未满足的先修数量，queue 保存当前可学习课程，visited\_count 是已处理数量。手算时按这个协议跟踪：0 出队后，遍历它的后继 1，indegree[1] 从 1 变 0，1 入队；最后处理数量等于课程数，说明无环。

\textbf{代码落点。} 关键是边方向要统一：prerequisite [a, b] 表示 b -> a。处理数量小于 numCourses 时返回 false。关键代码行必须能逐句翻译成“旧状态怎样变成新状态”，否则就只是模板。

\textbf{正确性。} 入度为 0 的节点没有未满足前置条件，可以安全学习；有向环中的节点永远不会全部降到入度 0。

\textbf{边界实验。} 先用边方向、孤立课程、环、自依赖、重复边做反例检查。实验时覆盖边方向、孤立课程、环、自依赖、重复边，并打印入队时的完整状态。若同一位置但资源不同的状态被合并，很多图题会出现隐蔽错误。

\textbf{复杂度变化。} 反复扫描最坏 O(VE)，拓扑排序 O(V+E) 时间和空间。

\textbf{迁移追问。} 如果题目改成“DFS 三色标记也能判环，但要明确访问状态”，先判断原状态是否仍然覆盖新答案集合，再决定是微调边界、改 value 语义，还是换算法。

\subsection{LeetCode 1334 阈值距离内邻居最少的城市}

\textbf{题目说明。} LeetCode 1334 阈值距离内邻居最少的城市 的题面入口要先还原成具体任务：每个城市都需要知道到其他城市的最短距离是否不超过阈值。

\textbf{样例入口。} 拿四个城市和距离阈值 4 手算，每个城市需要统计最短距离不超过 4 的其他城市数；并列时选编号更大的城市。

\textbf{暴力基线。} 慢版本可以枚举路径，或先用普通队列试跑一个小图。它会暴露步数最少和代价最小是否一致；一旦不一致，就必须围绕代价组织状态。放到本题就是：先按“每个城市都需要知道到其他城市的最短距离是否不超过阈值”完整试慢版本；样例里已经能看到“规律是先求全源最短路，再按邻居数量和编号规则选答案”，所以优化前必须先能说清慢版本枚举了哪些候选。

\textbf{暴力过程。} 拿四个城市和距离阈值 4 手算，每个城市需要统计最短距离不超过 4 的其他城市数；并列时选编号更大的城市。

\textbf{重复工作。} 题面模型：每个城市都需要知道到其他城市的最短距离是否不超过阈值。暴力版的慢点要落到本题：慢版本会围绕“每个城市都需要知道到其他城市的最短距离是否不超过阈值”枚举路径或按错误顺序扩展状态。样例规律是“规律是先求全源最短路，再按邻居数量和编号规则选答案”，说明队列或堆的弹出顺序必须和代价定义一致，旧状态为什么能跳过也要讲清。后面的优化只消掉这类重复工作，不能改变答案集合。

\textbf{优化条件。} 本题的关键观察是：节点数较小时可用 Floyd，边稀疏时也可从每个点跑 Dijkstra。这句话必须能翻译成可维护状态；如果题目条件改变后观察不再成立，就不能继续套原来的写法。

\textbf{相邻状态。} 规律是先求全源最短路，再按邻居数量和编号规则选答案。把这个特例继续拆开看：队列或堆弹出的顺序必须和代价规则一致；旧状态被跳过，是因为已经有更小代价到达同一状态。只要边权或代价合成方式改变，就要重新证明扩展顺序。

\textbf{状态复用。} 把路径枚举压成距离数组和优先队列；每次只扩展当前最有希望的未过期状态。本题落点是：规律是先求全源最短路，再按邻居数量和编号规则选答案。手算时只改被新输入影响的那一小部分，而不是回头重算完整候选。

\textbf{指针和字段。} 状态字段要能说清：距离数组、优先队列状态、松弛候选、旧状态判定、不可达哨兵；资源约束题还要把资源维度放进状态。手算时按这个协议跟踪：拿四个城市和距离阈值 4 手算，每个城市需要统计最短距离不超过 4 的其他城市数；并列时选编号更大的城市。然后按协议复查：手算时记录每次弹出的代价、当前距离数组和松弛结果。旧状态被弹出时为什么跳过，也要写在表里。

\textbf{代码落点。} 松弛时只在候选更优时更新；priority\_queue 弹出旧状态要跳过；距离加法用更宽整数。关键代码行必须能逐句翻译成“旧状态怎样变成新状态”，否则就只是模板。

\textbf{正确性。} 证明从扩展顺序开始：当前弹出的状态在代价规则下已经不会被未来路径改得更优。

\textbf{边界实验。} 先用距离等于阈值、并列时选编号最大、不可达、无向边做反例检查。实验时覆盖距离等于阈值、并列时选编号最大、不可达、无向边，打印每次弹出的代价和是否过期。若普通 BFS 与 Dijkstra 结果不同，要能解释边权条件造成的差异。

\textbf{复杂度变化。} 暴力枚举路径可能指数级；Dijkstra 在邻接表和堆实现下按边数、点数和堆操作计算，0-1 BFS 可按节点数加边数计算。

\textbf{迁移追问。} 如果题目改成“这是多源最短路选择问题，算法取决于 n、边数和查询密度”，先判断原状态是否仍然覆盖新答案集合，再决定是微调边界、改 value 语义，还是换算法。

\subsection{LeetCode 1584 连接所有点的最小费用}

\textbf{题目说明。} LeetCode 1584 连接所有点的最小费用 的题面入口要先还原成具体任务：完全图上连接所有点的最小总成本是最小生成树问题，边权为曼哈顿距离。

\textbf{样例入口。} 拿 points=[[0, 0], [2, 2], [3, 10], [5, 2], [7, 0]] 手算，任意两点都有曼哈顿距离边，目标是选 n-1 条边连通且总成本最小。

\textbf{暴力基线。} 慢版本先按题意画出完整状态图，用普通搜索跑通可达性。这个过程用于确认完全图上连接所有点的最小总成本是最小生成树问题，边权为曼哈顿距离的节点和边，之后再讨论访问标记、层次或代价顺序。放到本题就是：先按“完全图上连接所有点的最小总成本是最小生成树问题，边权为曼哈顿距离”完整试慢版本；样例里已经能看到“规律是最小生成树问题，Prim 维护每个外部点到已连通集合的最小边，每次加入最便宜外部点”，所以优化前必须先能说清慢版本枚举了哪些候选。

\textbf{暴力过程。} 拿 points=[[0, 0], [2, 2], [3, 10], [5, 2], [7, 0]] 手算，任意两点都有曼哈顿距离边，目标是选 n-1 条边连通且总成本最小。

\textbf{重复工作。} 题面模型：完全图上连接所有点的最小总成本是最小生成树问题，边权为曼哈顿距离。暴力版的慢点要落到本题：慢版本会围绕“完全图上连接所有点的最小总成本是最小生成树问题，边权为曼哈顿距离”反复展开同一批状态。样例规律是“规律是最小生成树问题，Prim 维护每个外部点到已连通集合的最小边，每次加入最便宜外部点”，说明必须把节点、边、访问标记和资源维度一次建清；同一状态不应重复入队，不同资源状态也不能误合并。后面的优化只消掉这类重复工作，不能改变答案集合。

\textbf{优化条件。} 本题的关键观察是：Prim 从已连通点集扩展，每次选到外部点的最小距离；也可 Kruskal 显式建边排序。这句话必须能翻译成可维护状态；如果题目条件改变后观察不再成立，就不能继续套原来的写法。

\textbf{相邻状态。} 规律是最小生成树问题，Prim 维护每个外部点到已连通集合的最小边，每次加入最便宜外部点。把这个特例继续拆开看：状态节点不一定等于原始位置，还可能包含钥匙、剩余资源、时间或访问集合。visited 只能合并未来能力完全相同的状态，否则会把合法路径剪掉。

\textbf{状态复用。} 把重复搜索压成访问标记、距离数组或拓扑状态；邻居生成也要从全量扫描变成结构化枚举。本题落点是：规律是最小生成树问题，Prim 维护每个外部点到已连通集合的最小边，每次加入最便宜外部点。手算时只改被新输入影响的那一小部分，而不是回头重算完整候选。

\textbf{指针和字段。} 状态字段要能说清：状态编号、邻接生成方式、访问标记、距离或层数、队列内容；若状态含资源，visited 不能只按位置记录。手算时按这个协议跟踪：拿 points=[[0, 0], [2, 2], [3, 10], [5, 2], [7, 0]] 手算，任意两点都有曼哈顿距离边，目标是选 n-1 条边连通且总成本最小。然后按协议复查：手算时画队列或栈，状态必须包含题目需要的所有资源。入队时标记还是出队时标记，要和去重语义一致。

\textbf{代码落点。} 入队时标记还是出队时标记必须一致；方向数组集中定义；多源 BFS 要一次性把所有源点放入初始队列。关键代码行必须能逐句翻译成“旧状态怎样变成新状态”，否则就只是模板。

\textbf{正确性。} 证明从可达性开始：搜索覆盖所有合法边能到达的状态，访问标记只删除重复状态而不删除不同未来能力。

\textbf{边界实验。} 先用单点、重复坐标、完全图边数、曼哈顿距离、Prim 更新外部点最小连边做反例检查。实验时覆盖单点、重复坐标、完全图边数、曼哈顿距离、Prim 更新外部点最小连边，并打印入队时的完整状态。若同一位置但资源不同的状态被合并，很多图题会出现隐蔽错误。

\textbf{复杂度变化。} 暴力重复从多个起点搜索会反复遍历图；正确建模和访问标记后，BFS 或 DFS 通常按节点数加边数计算，状态图题则按完整状态数计算。

\textbf{迁移追问。} 如果题目改成“规模很大时要研究几何图优化，不能显式 O(\ensuremath{n^{2}}) 建边”，先判断原状态是否仍然覆盖新答案集合，再决定是微调边界、改 value 语义，还是换算法。

\subsection{LeetCode 1192 查找集群内的关键连接}

\textbf{题目说明。} LeetCode 1192 查找集群内的关键连接 的题面入口要先还原成具体任务：无向图中的桥是删除后会断开连通性的边。

\textbf{样例入口。} 拿 n=4、connections=[[0, 1], [1, 2], [2, 0], [1, 3]] 手算，0-1-2 形成环，删除其中任一条仍可互通；节点 3 只通过边 1-3 接入，删除它会断开。

\textbf{暴力基线。} 慢版本先按题意画出完整状态图，用普通搜索跑通可达性。这个过程用于确认无向图中的桥是删除后会断开连通性的边的节点和边，之后再讨论访问标记、层次或代价顺序。放到本题就是：先按“无向图中的桥是删除后会断开连通性的边”完整试慢版本；样例里已经能看到“规律是 DFS 中 low[child] 若大于 disc[parent]，说明 child 子树没有备用返祖通道，这条边就是桥”，所以优化前必须先能说清慢版本枚举了哪些候选。

\textbf{暴力过程。} 拿 n=4、connections=[[0, 1], [1, 2], [2, 0], [1, 3]] 手算，0-1-2 形成环，删除其中任一条仍可互通；节点 3 只通过边 1-3 接入，删除它会断开。

\textbf{重复工作。} 题面模型：无向图中的桥是删除后会断开连通性的边。暴力版的慢点要落到本题：慢版本会围绕“无向图中的桥是删除后会断开连通性的边”反复展开同一批状态。样例规律是“规律是 DFS 中 low[child] 若大于 disc[parent]，说明 child 子树没有备用返祖通道，这条边就是桥”，说明必须把节点、边、访问标记和资源维度一次建清；同一状态不应重复入队，不同资源状态也不能误合并。后面的优化只消掉这类重复工作，不能改变答案集合。

\textbf{优化条件。} 本题的关键观察是：DFS 记录 disc 和 low，若 low[child] 大于 disc[parent]，子树没有返祖通道，父子边是桥。这句话必须能翻译成可维护状态；如果题目条件改变后观察不再成立，就不能继续套原来的写法。

\textbf{相邻状态。} 规律是 DFS 中 low[child] 若大于 disc[parent]，说明 child 子树没有备用返祖通道，这条边就是桥。把这个特例继续拆开看：状态节点不一定等于原始位置，还可能包含钥匙、剩余资源、时间或访问集合。visited 只能合并未来能力完全相同的状态，否则会把合法路径剪掉。

\textbf{状态复用。} 把重复搜索压成访问标记、距离数组或拓扑状态；邻居生成也要从全量扫描变成结构化枚举。本题落点是：规律是 DFS 中 low[child] 若大于 disc[parent]，说明 child 子树没有备用返祖通道，这条边就是桥。手算时只改被新输入影响的那一小部分，而不是回头重算完整候选。

\textbf{指针和字段。} 状态字段要能说清：状态编号、邻接生成方式、访问标记、距离或层数、队列内容；若状态含资源，visited 不能只按位置记录。手算时按这个协议跟踪：拿 n=4、connections=[[0, 1], [1, 2], [2, 0], [1, 3]] 手算，0-1-2 形成环，删除其中任一条仍可互通；节点 3 只通过边 1-3 接入，删除它会断开。然后按协议复查：手算时画队列或栈，状态必须包含题目需要的所有资源。入队时标记还是出队时标记，要和去重语义一致。

\textbf{代码落点。} 入队时标记还是出队时标记必须一致；方向数组集中定义；多源 BFS 要一次性把所有源点放入初始队列。关键代码行必须能逐句翻译成“旧状态怎样变成新状态”，否则就只是模板。

\textbf{正确性。} 证明从可达性开始：搜索覆盖所有合法边能到达的状态，访问标记只删除重复状态而不删除不同未来能力。

\textbf{边界实验。} 先用重边、根节点、图保证连通、无向边需要编号跳过反向边、递归栈做反例检查。实验时覆盖重边、根节点、图保证连通、无向边需要编号跳过反向边、递归栈，并打印入队时的完整状态。若同一位置但资源不同的状态被合并，很多图题会出现隐蔽错误。

\textbf{复杂度变化。} 暴力重复从多个起点搜索会反复遍历图；正确建模和访问标记后，BFS 或 DFS 通常按节点数加边数计算，状态图题则按完整状态数计算。

\textbf{迁移追问。} 如果题目改成“若要找割点，判断条件和根节点处理不同”，先判断原状态是否仍然覆盖新答案集合，再决定是微调边界、改 value 语义，还是换算法。

\subsection{LeetCode 1820 最多邀请的个数}

\textbf{题目说明。} LeetCode 1820 最多邀请的个数 的题面入口要先还原成具体任务：男生和女生构成二分图，目标是最大一对一匹配数量。

\textbf{样例入口。} 拿 grid=[[1, 1, 1], [1, 0, 1], [0, 0, 1]] 手算，若第一个男生先匹配第一个女生，后面男生仍可能通过重新分配得到三对。贪心固定早期选择会挡住后续唯一选择。

\textbf{暴力基线。} 慢版本先按题意画出完整状态图，用普通搜索跑通可达性。这个过程用于确认男生和女生构成二分图，目标是最大一对一匹配数量的节点和边，之后再讨论访问标记、层次或代价顺序。放到本题就是：先按“男生和女生构成二分图，目标是最大一对一匹配数量”完整试慢版本；样例里已经能看到“规律是每次寻找增广路，找到未匹配右点后翻转路径上的匹配关系，使匹配数增加一”，所以优化前必须先能说清慢版本枚举了哪些候选。

\textbf{暴力过程。} 拿 grid=[[1, 1, 1], [1, 0, 1], [0, 0, 1]] 手算，若第一个男生先匹配第一个女生，后面男生仍可能通过重新分配得到三对。贪心固定早期选择会挡住后续唯一选择。

\textbf{重复工作。} 题面模型：男生和女生构成二分图，目标是最大一对一匹配数量。暴力版的慢点要落到本题：慢版本会围绕“男生和女生构成二分图，目标是最大一对一匹配数量”反复展开同一批状态。样例规律是“规律是每次寻找增广路，找到未匹配右点后翻转路径上的匹配关系，使匹配数增加一”，说明必须把节点、边、访问标记和资源维度一次建清；同一状态不应重复入队，不同资源状态也不能误合并。后面的优化只消掉这类重复工作，不能改变答案集合。

\textbf{优化条件。} 本题的关键观察是：对每个未匹配左点寻找增广路，找到未匹配右点后翻转路径，使匹配数加一。这句话必须能翻译成可维护状态；如果题目条件改变后观察不再成立，就不能继续套原来的写法。

\textbf{相邻状态。} 规律是每次寻找增广路，找到未匹配右点后翻转路径上的匹配关系，使匹配数增加一。把这个特例继续拆开看：状态节点不一定等于原始位置，还可能包含钥匙、剩余资源、时间或访问集合。visited 只能合并未来能力完全相同的状态，否则会把合法路径剪掉。

\textbf{状态复用。} 把重复搜索压成访问标记、距离数组或拓扑状态；邻居生成也要从全量扫描变成结构化枚举。本题落点是：规律是每次寻找增广路，找到未匹配右点后翻转路径上的匹配关系，使匹配数增加一。手算时只改被新输入影响的那一小部分，而不是回头重算完整候选。

\textbf{指针和字段。} 状态字段要能说清：状态编号、邻接生成方式、访问标记、距离或层数、队列内容；若状态含资源，visited 不能只按位置记录。手算时按这个协议跟踪：拿 grid=[[1, 1, 1], [1, 0, 1], [0, 0, 1]] 手算，若第一个男生先匹配第一个女生，后面男生仍可能通过重新分配得到三对。贪心固定早期选择会挡住后续唯一选择。然后按协议复查：手算时画队列或栈，状态必须包含题目需要的所有资源。入队时标记还是出队时标记，要和去重语义一致。

\textbf{代码落点。} 入队时标记还是出队时标记必须一致；方向数组集中定义；多源 BFS 要一次性把所有源点放入初始队列。关键代码行必须能逐句翻译成“旧状态怎样变成新状态”，否则就只是模板。

\textbf{正确性。} 证明从可达性开始：搜索覆盖所有合法边能到达的状态，访问标记只删除重复状态而不删除不同未来能力。

\textbf{边界实验。} 先用某一侧为空、多个男生争同一女生、贪心失败、访问标记按轮重置、稀疏边做反例检查。实验时覆盖某一侧为空、多个男生争同一女生、贪心失败、访问标记按轮重置、稀疏边，并打印入队时的完整状态。若同一位置但资源不同的状态被合并，很多图题会出现隐蔽错误。

\textbf{复杂度变化。} 暴力重复从多个起点搜索会反复遍历图；正确建模和访问标记后，BFS 或 DFS 通常按节点数加边数计算，状态图题则按完整状态数计算。

\textbf{迁移追问。} 如果题目改成“若边有容量或权重，升级为最大流或带权匹配”，先判断原状态是否仍然覆盖新答案集合，再决定是微调边界、改 value 语义，还是换算法。

\subsection{LeetCode 721 账户合并}

\textbf{题目说明。} LeetCode 721 账户合并 的题面入口要先还原成具体任务：共享邮箱说明账户属于同一人，邮箱之间形成连通分量。

\textbf{样例入口。} 拿 John 的账号 [a, b] 和另一个 [b, c] 手算，邮箱 b 重叠，所以 a、b、c 属于同一人；名字只是输出标签，合并依据是邮箱。

\textbf{暴力基线。} 慢版本每次合并后用 DFS 或 BFS 重新判断连通性。它正确但重复遍历已有连接；当关系只会合并不会拆开时，可以压成集合代表。放到本题就是：先按“共享邮箱说明账户属于同一人，邮箱之间形成连通分量”完整试慢版本；样例里已经能看到“规律是邮箱映射到编号后并查集合并，最后按代表收集并排序邮箱”，所以优化前必须先能说清慢版本枚举了哪些候选。

\textbf{暴力过程。} 拿 John 的账号 [a, b] 和另一个 [b, c] 手算，邮箱 b 重叠，所以 a、b、c 属于同一人；名字只是输出标签，合并依据是邮箱。

\textbf{重复工作。} 题面模型：共享邮箱说明账户属于同一人，邮箱之间形成连通分量。暴力版的慢点要落到本题：慢版本会围绕“共享邮箱说明账户属于同一人，邮箱之间形成连通分量”反复搜索连通性。样例规律是“规律是邮箱映射到编号后并查集合并，最后按代表收集并排序邮箱”，说明关系只需要被压成集合代表；每次 union 失败或成功都要能翻译回题意。后面的优化只消掉这类重复工作，不能改变答案集合。

\textbf{优化条件。} 本题的关键观察是：给邮箱编号并查集合并，同根邮箱收集后排序输出。这句话必须能翻译成可维护状态；如果题目条件改变后观察不再成立，就不能继续套原来的写法。

\textbf{相邻状态。} 规律是邮箱映射到编号后并查集合并，最后按代表收集并排序邮箱。把这个特例继续拆开看：union 只是在维护等价类，不是在保存具体路径。两个节点代表元相同只能说明连通或关系一致；如果题目还要距离、比例或奇偶性，必须把附加信息挂到根关系上。

\textbf{状态复用。} 把连通分量压成代表元；查询只比较代表，合并只发生在两个根之间。本题落点是：规律是邮箱映射到编号后并查集合并，最后按代表收集并排序邮箱。手算时只改被新输入影响的那一小部分，而不是回头重算完整候选。

\textbf{指针和字段。} 状态字段要能说清：parent、size 或 rank、find 返回的代表元、合并时的附加信息；带权或奇偶并查集还要维护到根的关系。手算时按这个协议跟踪：拿 John 的账号 [a, b] 和另一个 [b, c] 手算，邮箱 b 重叠，所以 a、b、c 属于同一人；名字只是输出标签，合并依据是邮箱。然后按协议复查：手算时写 parent 表和集合代表。每次 union 只改变根之间的关系，find 后代表相同才说明连通。

\textbf{代码落点。} find 路径压缩不能破坏附加权值；union 只能连接两个根；合并失败和重复边要保留题意解释。关键代码行必须能逐句翻译成“旧状态怎样变成新状态”，否则就只是模板。

\textbf{正确性。} 证明从等价关系开始：合并维护连通性的传递性，find 返回的代表元就是当前集合身份。

\textbf{边界实验。} 先用同名不同人、一个账户多个邮箱、重复邮箱、输出排序做反例检查。实验时覆盖同名不同人、一个账户多个邮箱、重复邮箱、输出排序，打印每个元素的代表元和集合大小。重复合并、已经连通的边和字符串编号最容易暴露实现问题。

\textbf{复杂度变化。} 暴力每次查询都搜索连通性代价高；并查集把合并和查询摊还到近似常数，整体接近线性。

\textbf{迁移追问。} 如果题目改成“姓名不能作为合并依据，只能作为输出字段”，先判断原状态是否仍然覆盖新答案集合，再决定是微调边界、改 value 语义，还是换算法。

\subsection{LeetCode 785 判断二分图}

\textbf{题目说明。} LeetCode 785 判断二分图 的题面入口要先还原成具体任务：二分图要求每条边两端颜色相反。

\textbf{样例入口。} 拿 graph=[[1, 2, 3], [0, 2], [0, 1, 3], [0, 2]] 手算，0 染红后 1、2、3 都要染蓝，但 1 和 2 之间还有边，蓝蓝相邻冲突，所以不是二分图。

\textbf{暴力基线。} 慢版本先按题意画出完整状态图，用普通搜索跑通可达性。这个过程用于确认二分图要求每条边两端颜色相反的节点和边，之后再讨论访问标记、层次或代价顺序。放到本题就是：先按“二分图要求每条边两端颜色相反”完整试慢版本；样例里已经能看到“规律是 BFS 染色时每条边两端必须异色，发现同色边立即失败”，所以优化前必须先能说清慢版本枚举了哪些候选。

\textbf{暴力过程。} 拿 graph=[[1, 2, 3], [0, 2], [0, 1, 3], [0, 2]] 手算，0 染红后 1、2、3 都要染蓝，但 1 和 2 之间还有边，蓝蓝相邻冲突，所以不是二分图。

\textbf{重复工作。} 题面模型：二分图要求每条边两端颜色相反。暴力版的慢点要落到本题：慢版本会围绕“二分图要求每条边两端颜色相反”反复展开同一批状态。样例规律是“规律是 BFS 染色时每条边两端必须异色，发现同色边立即失败”，说明必须把节点、边、访问标记和资源维度一次建清；同一状态不应重复入队，不同资源状态也不能误合并。后面的优化只消掉这类重复工作，不能改变答案集合。

\textbf{优化条件。} 本题的关键观察是：对每个未染色连通块 BFS 或 DFS 染色，遇到同色边立即失败。这句话必须能翻译成可维护状态；如果题目条件改变后观察不再成立，就不能继续套原来的写法。

\textbf{相邻状态。} 规律是 BFS 染色时每条边两端必须异色，发现同色边立即失败。把这个特例继续拆开看：状态节点不一定等于原始位置，还可能包含钥匙、剩余资源、时间或访问集合。visited 只能合并未来能力完全相同的状态，否则会把合法路径剪掉。

\textbf{状态复用。} 把重复搜索压成访问标记、距离数组或拓扑状态；邻居生成也要从全量扫描变成结构化枚举。本题落点是：规律是 BFS 染色时每条边两端必须异色，发现同色边立即失败。手算时只改被新输入影响的那一小部分，而不是回头重算完整候选。

\textbf{指针和字段。} 状态字段要能说清：状态编号、邻接生成方式、访问标记、距离或层数、队列内容；若状态含资源，visited 不能只按位置记录。手算时按这个协议跟踪：拿 graph=[[1, 2, 3], [0, 2], [0, 1, 3], [0, 2]] 手算，0 染红后 1、2、3 都要染蓝，但 1 和 2 之间还有边，蓝蓝相邻冲突，所以不是二分图。然后按协议复查：手算时画队列或栈，状态必须包含题目需要的所有资源。入队时标记还是出队时标记，要和去重语义一致。

\textbf{代码落点。} 入队时标记还是出队时标记必须一致；方向数组集中定义；多源 BFS 要一次性把所有源点放入初始队列。关键代码行必须能逐句翻译成“旧状态怎样变成新状态”，否则就只是模板。

\textbf{正确性。} 证明从可达性开始：搜索覆盖所有合法边能到达的状态，访问标记只删除重复状态而不删除不同未来能力。

\textbf{边界实验。} 先用非连通图、自环、奇环、孤立点、重复边做反例检查。实验时覆盖非连通图、自环、奇环、孤立点、重复边，并打印入队时的完整状态。若同一位置但资源不同的状态被合并，很多图题会出现隐蔽错误。

\textbf{复杂度变化。} 暴力重复从多个起点搜索会反复遍历图；正确建模和访问标记后，BFS 或 DFS 通常按节点数加边数计算，状态图题则按完整状态数计算。

\textbf{迁移追问。} 如果题目改成“可能的二分法是同一染色模型，只是图由 dislike 关系构造”，先判断原状态是否仍然覆盖新答案集合，再决定是微调边界、改 value 语义，还是换算法。

\subsection{LeetCode 752 打开转盘锁}

\textbf{题目说明。} LeetCode 752 打开转盘锁 的题面入口要先还原成具体任务：四位密码是状态，每次转动一位形成边。

\textbf{样例入口。} 拿锁 0000 到 0001 手算，每一步只能拨动一位加一或减一，0000 的邻居有 1000、9000、0100 等。

\textbf{暴力基线。} 慢版本先按题意画出完整状态图，用普通搜索跑通可达性。这个过程用于确认四位密码是状态，每次转动一位形成边的节点和边，之后再讨论访问标记、层次或代价顺序。放到本题就是：先按“四位密码是状态，每次转动一位形成边”完整试慢版本；样例里已经能看到“规律是状态是四位字符串，BFS 按步数扩展，deadend 入队前就要排除”，所以优化前必须先能说清慢版本枚举了哪些候选。

\textbf{暴力过程。} 拿锁 0000 到 0001 手算，每一步只能拨动一位加一或减一，0000 的邻居有 1000、9000、0100 等。

\textbf{重复工作。} 题面模型：四位密码是状态，每次转动一位形成边。暴力版的慢点要落到本题：慢版本会围绕“四位密码是状态，每次转动一位形成边”反复展开同一批状态。样例规律是“规律是状态是四位字符串，BFS 按步数扩展，deadend 入队前就要排除”，说明必须把节点、边、访问标记和资源维度一次建清；同一状态不应重复入队，不同资源状态也不能误合并。后面的优化只消掉这类重复工作，不能改变答案集合。

\textbf{优化条件。} 本题的关键观察是：BFS 从起点扩展，死亡状态不能入队，访问状态避免重复。这句话必须能翻译成可维护状态；如果题目条件改变后观察不再成立，就不能继续套原来的写法。

\textbf{相邻状态。} 规律是状态是四位字符串，BFS 按步数扩展，deadend 入队前就要排除。把这个特例继续拆开看：状态节点不一定等于原始位置，还可能包含钥匙、剩余资源、时间或访问集合。visited 只能合并未来能力完全相同的状态，否则会把合法路径剪掉。

\textbf{状态复用。} 把重复搜索压成访问标记、距离数组或拓扑状态；邻居生成也要从全量扫描变成结构化枚举。本题落点是：规律是状态是四位字符串，BFS 按步数扩展，deadend 入队前就要排除。手算时只改被新输入影响的那一小部分，而不是回头重算完整候选。

\textbf{指针和字段。} 状态字段要能说清：状态编号、邻接生成方式、访问标记、距离或层数、队列内容；若状态含资源，visited 不能只按位置记录。手算时按这个协议跟踪：拿锁 0000 到 0001 手算，每一步只能拨动一位加一或减一，0000 的邻居有 1000、9000、0100 等。然后按协议复查：手算时画队列或栈，状态必须包含题目需要的所有资源。入队时标记还是出队时标记，要和去重语义一致。

\textbf{代码落点。} 入队时标记还是出队时标记必须一致；方向数组集中定义；多源 BFS 要一次性把所有源点放入初始队列。关键代码行必须能逐句翻译成“旧状态怎样变成新状态”，否则就只是模板。

\textbf{正确性。} 证明从可达性开始：搜索覆盖所有合法边能到达的状态，访问标记只删除重复状态而不删除不同未来能力。

\textbf{边界实验。} 先用起点死亡、目标就是起点、数字环绕、死亡集合很大做反例检查。实验时覆盖起点死亡、目标就是起点、数字环绕、死亡集合很大，并打印入队时的完整状态。若同一位置但资源不同的状态被合并，很多图题会出现隐蔽错误。

\textbf{复杂度变化。} 暴力重复从多个起点搜索会反复遍历图；正确建模和访问标记后，BFS 或 DFS 通常按节点数加边数计算，状态图题则按完整状态数计算。

\textbf{迁移追问。} 如果题目改成“双向 BFS 能明显减少状态展开”，先判断原状态是否仍然覆盖新答案集合，再决定是微调边界、改 value 语义，还是换算法。

\subsection{LeetCode 1293 网格中的最短路径}

\textbf{题目说明。} LeetCode 1293 网格中的最短路径 的题面入口要先还原成具体任务：走到同一格时，剩余可消除障碍数量不同会影响未来可达性。

\textbf{样例入口。} 拿网格三列中间有障碍且 k=1 手算，走到同一格时，如果剩余消障次数不同，未来可走路径也不同。

\textbf{暴力基线。} 慢版本先按题意画出完整状态图，用普通搜索跑通可达性。这个过程用于确认走到同一格时，剩余可消除障碍数量不同会影响未来可达性的节点和边，之后再讨论访问标记、层次或代价顺序。放到本题就是：先按“走到同一格时，剩余可消除障碍数量不同会影响未来可达性”完整试慢版本；样例里已经能看到“规律是 BFS 的状态必须包含剩余 k，visited 不能只按位置记录”，所以优化前必须先能说清慢版本枚举了哪些候选。

\textbf{暴力过程。} 拿网格三列中间有障碍且 k=1 手算，走到同一格时，如果剩余消障次数不同，未来可走路径也不同。

\textbf{重复工作。} 题面模型：走到同一格时，剩余可消除障碍数量不同会影响未来可达性。暴力版的慢点要落到本题：慢版本会围绕“走到同一格时，剩余可消除障碍数量不同会影响未来可达性”反复展开同一批状态。样例规律是“规律是 BFS 的状态必须包含剩余 k，visited 不能只按位置记录”，说明必须把节点、边、访问标记和资源维度一次建清；同一状态不应重复入队，不同资源状态也不能误合并。后面的优化只消掉这类重复工作，不能改变答案集合。

\textbf{优化条件。} 本题的关键观察是：BFS 状态包含位置和剩余消除次数，visited 记录该维度或记录到达该格的最大剩余次数。这句话必须能翻译成可维护状态；如果题目条件改变后观察不再成立，就不能继续套原来的写法。

\textbf{相邻状态。} 规律是 BFS 的状态必须包含剩余 k，visited 不能只按位置记录。把这个特例继续拆开看：状态节点不一定等于原始位置，还可能包含钥匙、剩余资源、时间或访问集合。visited 只能合并未来能力完全相同的状态，否则会把合法路径剪掉。

\textbf{状态复用。} 把重复搜索压成访问标记、距离数组或拓扑状态；邻居生成也要从全量扫描变成结构化枚举。本题落点是：规律是 BFS 的状态必须包含剩余 k，visited 不能只按位置记录。手算时只改被新输入影响的那一小部分，而不是回头重算完整候选。

\textbf{指针和字段。} 状态字段要能说清：状态编号、邻接生成方式、访问标记、距离或层数、队列内容；若状态含资源，visited 不能只按位置记录。手算时按这个协议跟踪：拿网格三列中间有障碍且 k=1 手算，走到同一格时，如果剩余消障次数不同，未来可走路径也不同。然后按协议复查：手算时画队列或栈，状态必须包含题目需要的所有资源。入队时标记还是出队时标记，要和去重语义一致。

\textbf{代码落点。} 入队时标记还是出队时标记必须一致；方向数组集中定义；多源 BFS 要一次性把所有源点放入初始队列。关键代码行必须能逐句翻译成“旧状态怎样变成新状态”，否则就只是模板。

\textbf{正确性。} 证明从可达性开始：搜索覆盖所有合法边能到达的状态，访问标记只删除重复状态而不删除不同未来能力。

\textbf{边界实验。} 先用起点终点、k 很大、障碍密集、二维 visited 错误剪枝做反例检查。实验时覆盖起点终点、k 很大、障碍密集、二维 visited 错误剪枝，并打印入队时的完整状态。若同一位置但资源不同的状态被合并，很多图题会出现隐蔽错误。

\textbf{复杂度变化。} 暴力重复从多个起点搜索会反复遍历图；正确建模和访问标记后，BFS 或 DFS 通常按节点数加边数计算，状态图题则按完整状态数计算。

\textbf{迁移追问。} 如果题目改成“它和钥匙题一样说明状态维度不能丢”，先判断原状态是否仍然覆盖新答案集合，再决定是微调边界、改 value 语义，还是换算法。

\subsection{LeetCode 864 获取所有钥匙的最短路径}

\textbf{题目说明。} LeetCode 864 获取所有钥匙的最短路径 的题面入口要先还原成具体任务：位置相同但钥匙集合不同，未来能开的门不同。

\textbf{样例入口。} 拿网格 @.a / \#.\# / b.A 手算，走到 a 后状态不是位置变了而是钥匙集合多了一把；同一位置带不同钥匙未来能力不同。

\textbf{暴力基线。} 慢版本先按题意画出完整状态图，用普通搜索跑通可达性。这个过程用于确认位置相同但钥匙集合不同，未来能开的门不同的节点和边，之后再讨论访问标记、层次或代价顺序。放到本题就是：先按“位置相同但钥匙集合不同，未来能开的门不同”完整试慢版本；样例里已经能看到“规律是 BFS 的 visited 必须按 (row, col, keymask) 记录，拿齐所有钥匙第一次出队就是最短步数”，所以优化前必须先能说清慢版本枚举了哪些候选。

\textbf{暴力过程。} 拿网格 @.a / \#.\# / b.A 手算，走到 a 后状态不是位置变了而是钥匙集合多了一把；同一位置带不同钥匙未来能力不同。

\textbf{重复工作。} 题面模型：位置相同但钥匙集合不同，未来能开的门不同。暴力版的慢点要落到本题：慢版本会围绕“位置相同但钥匙集合不同，未来能开的门不同”反复展开同一批状态。样例规律是“规律是 BFS 的 visited 必须按 (row, col, keymask) 记录，拿齐所有钥匙第一次出队就是最短步数”，说明必须把节点、边、访问标记和资源维度一次建清；同一状态不应重复入队，不同资源状态也不能误合并。后面的优化只消掉这类重复工作，不能改变答案集合。

\textbf{优化条件。} 本题的关键观察是：BFS 状态包含行、列和钥匙掩码，访问标记也必须包含掩码。这句话必须能翻译成可维护状态；如果题目条件改变后观察不再成立，就不能继续套原来的写法。

\textbf{相邻状态。} 规律是 BFS 的 visited 必须按 (row, col, keymask) 记录，拿齐所有钥匙第一次出队就是最短步数。把这个特例继续拆开看：状态节点不一定等于原始位置，还可能包含钥匙、剩余资源、时间或访问集合。visited 只能合并未来能力完全相同的状态，否则会把合法路径剪掉。

\textbf{状态复用。} 把重复搜索压成访问标记、距离数组或拓扑状态；邻居生成也要从全量扫描变成结构化枚举。本题落点是：规律是 BFS 的 visited 必须按 (row, col, keymask) 记录，拿齐所有钥匙第一次出队就是最短步数。手算时只改被新输入影响的那一小部分，而不是回头重算完整候选。

\textbf{指针和字段。} 状态字段要能说清：状态编号、邻接生成方式、访问标记、距离或层数、队列内容；若状态含资源，visited 不能只按位置记录。手算时按这个协议跟踪：拿网格 @.a / \#.\# / b.A 手算，走到 a 后状态不是位置变了而是钥匙集合多了一把；同一位置带不同钥匙未来能力不同。然后按协议复查：手算时画队列或栈，状态必须包含题目需要的所有资源。入队时标记还是出队时标记，要和去重语义一致。

\textbf{代码落点。} 入队时标记还是出队时标记必须一致；方向数组集中定义；多源 BFS 要一次性把所有源点放入初始队列。关键代码行必须能逐句翻译成“旧状态怎样变成新状态”，否则就只是模板。

\textbf{正确性。} 证明从可达性开始：搜索覆盖所有合法边能到达的状态，访问标记只删除重复状态而不删除不同未来能力。

\textbf{边界实验。} 先用起点旁边有门、钥匙顺序、不可达、钥匙数量决定掩码大小做反例检查。实验时覆盖起点旁边有门、钥匙顺序、不可达、钥匙数量决定掩码大小，并打印入队时的完整状态。若同一位置但资源不同的状态被合并，很多图题会出现隐蔽错误。

\textbf{复杂度变化。} 暴力重复从多个起点搜索会反复遍历图；正确建模和访问标记后，BFS 或 DFS 通常按节点数加边数计算，状态图题则按完整状态数计算。

\textbf{迁移追问。} 如果题目改成“这是状态图维度不能只按位置去重的典型题”，先判断原状态是否仍然覆盖新答案集合，再决定是微调边界、改 value 语义，还是换算法。

\subsection{LeetCode 200 岛屿数量}

\textbf{题目说明。} LeetCode 200 岛屿数量 的题面入口要先还原成具体任务：陆地格子是节点，四方向相邻是边，岛屿是连通块。

\textbf{样例入口。} 拿一个 2x2 小岛手算，从一个陆地出发能沿上下左右标记同一块岛。若不标记访问，会重复计数；若对角线也连通，会误合并。

\textbf{暴力基线。} 暴力可以对每个陆地格子单独搜索它所属的连通块，再用集合去重；这样同一块岛会被重复搜索很多次。

\textbf{暴力过程。} 以 2x2 全是陆地为例，从左上搜索会访问四个格子；如果之后从右上又重新搜索，就把同一座岛重复数了一遍。

\textbf{重复工作。} 题面模型：陆地格子是节点，四方向相邻是边，岛屿是连通块。暴力版的慢点要落到本题：浪费在已确认属于某个岛的陆地没有被标记。每块岛只需要在第一次遇到时完整淹没或标记。后面的优化只消掉这类重复工作，不能改变答案集合。

\textbf{优化条件。} 本题的关键观察是：扫描遇到未访问陆地就启动搜索并标记整块。这句话必须能翻译成可维护状态；如果题目条件改变后观察不再成立，就不能继续套原来的写法。

\textbf{相邻状态。} 规律是每发现一个未访问陆地，答案加一，并用 DFS/BFS 标记与它四连通的全部陆地。把这个特例继续拆开看：状态节点不一定等于原始位置，还可能包含钥匙、剩余资源、时间或访问集合。visited 只能合并未来能力完全相同的状态，否则会把合法路径剪掉。

\textbf{状态复用。} 扫描网格。遇到未访问陆地，答案加一，然后 BFS/DFS 把与它四方向连通的全部陆地标记为已访问。手算时只改被新输入影响的那一小部分，而不是回头重算完整候选。

\textbf{指针和字段。} 状态字段要能说清：row/col 是当前格子，visited 或原地改写表示已经归属某座岛，directions 是四方向邻居。手算时按这个协议跟踪：从左上陆地入队，依次扩展右边和下边陆地，最终整块连通区域都标记；扫描到这些格子时直接跳过。

\textbf{代码落点。} 关键是只看上下左右，不看对角线。边界判断集中写，访问标记在入队时设置，避免重复入队。关键代码行必须能逐句翻译成“旧状态怎样变成新状态”，否则就只是模板。

\textbf{正确性。} 一次搜索会访问且只访问与起点四连通的全部陆地；外层每次启动搜索都对应一块此前未计数的新岛。

\textbf{边界实验。} 先用空网格、全水、全陆、斜对角不连通、是否允许修改输入做反例检查。实验时覆盖空网格、全水、全陆、斜对角不连通、是否允许修改输入，并打印入队时的完整状态。若同一位置但资源不同的状态被合并，很多图题会出现隐蔽错误。

\textbf{复杂度变化。} 重复搜索最坏 O(\ensuremath{(mn)^{2}})，标记搜索 O(mn) 时间，空间 O(mn) 或递归/队列大小。

\textbf{迁移追问。} 如果题目改成“也可用并查集合并相邻陆地，适合动态岛屿扩展”，先判断原状态是否仍然覆盖新答案集合，再决定是微调边界、改 value 语义，还是换算法。

\begin{principlebox}{本节复盘方式}
不要只看最终算法名。每道题复盘时先遮住优化版，口述暴力枚举对象；再指出相邻窗口、历史前缀、候选区间、搜索状态或子问题里哪一部分被重复处理；最后用一个边界样例验证状态更新顺序。能讲完这三步，才算真正掌握本章案例。
\end{principlebox}
